\section{Discussion}
\label{sec:discussion}

\paragraph{Implications for practice.} Six suggestions follow from the tables, within the limits below.
\begin{enumerate}
  \item \emph{Studentize inside the bootstrap, or not at all.} A studentized data-snooping test whose resampled statistics reuse the sample's variance estimate is oversized at two years of daily data, more so with long blocks, many strategies or negative skew. Re-studentizing every resample keeps the test calibrated and keeps the power studentizing buys when volatilities differ. When strategies have similar volatilities, the Reality Check is a calibrated and simple alternative.
  \item \emph{Recentre consistently and count ties.} Hansen's consistent recentring makes the test far more powerful when many variants are clearly poor; computing the $p$-value as $\Pr^*(T^*\ge T)$ stops a floored statistic from rejecting when no strategy beats the benchmark.
  \item \emph{Check the selected strategy's returns for autocorrelation} before any multiple-testing adjustment. Lo's factor removes the excess in the autoregressive case it assumes; kernel and block-bootstrap corrections reduce it without that assumption, at a cost in size at this sample length.
  \item \emph{Compare corrections at matched sides and with the number of resamples stated.} Otherwise the comparison measures conventions: the haircut's conservatism comes from its two-sided convention, and a bootstrap's power from its number of resamples.
  \item \emph{Treat $\mathrm{DSR}\ge 0.95$ as the demanding rule it is.} It asks whether the best strategy beats the expected best of $N$ skill-less ones, not whether it has any skill, and over two years of daily data it passes few strategies with a true Sharpe ratio of 2.
  \item \emph{When returns are markedly skewed or fat-tailed, do not trust a single-series correction}: none in the benchmark is calibrated across negative skew, positive skew and fat tails, in part because selection itself distorts the moments of the chosen series. The Reality Check and the bootstrap-$t$ SPA test, which resample all strategies together, were.
\end{enumerate}

\paragraph{Limitations.} The nine families change one feature at a time; real returns combine several, and the only serial dependence studied is AR(1), the Lo factor's best case. The main grid covers one search shape ($N=\nz{trials}$, $T=\nz{obs}$) and one strongly skilled strategy; experiment B varies $T$ and $N$ under two families only, and real searches often contain many weakly skilled strategies. The block length of every resampling test is fixed by rule rather than chosen from the data \citep{politis2004automatic}. The bootstrap-$t$ studentizes each resample by its own standard deviation, which ignores serial dependence; a block-based studentization \citep{gotze1996second} would be the natural choice for autocorrelated returns, and is untested here. The deflated Sharpe ratio is scored only at its 0.95 point, not as a probability. Experiment C uses i.i.d.\ normal returns, so it isolates the role of studentizing and recentring but not their interaction with skew. Finally, a benchmark designed by the author of one of the scored methods, whose runs also chose a default for the author's own software, cannot be neutral in the sense of \citet{bartos2025living}. We pre-register the runs and publish every generator, validator, seed and result file, but independent replication with independently chosen families would be more convincing than anything we can do ourselves.

\paragraph{Next versions.} The benchmark is versioned so that it can grow without rewriting its history. Natural additions are families that combine features (skewed and autocorrelated returns) or have non-AR dependence, alternatives with many weakly skilled strategies, data-driven block lengths, block-studentized bootstrap tests, and an estimate of the effective number of trials for correlated searches.

\section{Conclusion}
\label{sec:conclusion}
The corrections practitioners use for backtest overfitting and data snooping keep their promises in the setting they were derived for and fail in predictable ways outside it. The most consequential failure found here is the most technical: Hansen's studentized SPA test, recommended over the Reality Check for its power, is oversized at the two-year horizon of strategy research because its bootstrap reuses a noisy variance estimate, and the arch implementation stays close to its level only because it does not do what its option says. Both are easy to repair. Autocorrelation remains the most damaging violation for the single-series corrections, and several apparent differences between methods come from conventions---sides, thresholds, numbers of resamples---rather than from the methods. An open, pre-registered benchmark with known ground truth makes such statements checkable, and anyone can rerun or extend it from the published code and seeds.
